\section{What these trajectories imply for agent evaluation}
\label{sec:discussion}

\subsection{Capability is a path through opportunities, not a single score}

The observed games contain economically consequential planning. Agents assemble assets through different routes, refuse exchanges that would finance rivals, mortgage spare collateral to build, retain reserves, and execute long liquidation sequences. They also miss available finance, submit recoverable allocations in invalid order, and receive actions that contradict their accompanying explanations. These behaviors coexist within configurations and sometimes within one trajectory. A win count compresses them into an endpoint without revealing which capabilities were exercised or which failures occurred.

The distinction between whole trajectories and conditioned opportunities is therefore substantive. The 43/63 first-window development count retains the effects of how a group was acquired and financed. The 34/42 entry-buildable count asks a narrower responsiveness question after that access already exists. Both are useful, and neither replaces the other. Conditioning away survival, management access or prior financing can remove part of the capability one intended to evaluate. Conversely, counting every late inherited group as an unused investment invents opportunities after the game has ended. A useful capability profile needs acquisition, access, applied execution and downstream consequences together.

The same issue arises in negotiation. Initiating few proposals can coexist with effective recipient bargaining; initiating many can reflect persistent but unsuccessful search. A complete group can come with insufficient cash, expensive redemption, a funded opponent or no future management. Two no-trade winners show that bargaining is not a necessary route in every realized game. The point is not to replace wins with another universal economic score, but to expose the decisions and resource transitions that make the outcome interpretable.

\subsection{Alignment needs a norm, an opportunity and an observable response}

The motivating question concerns both agent capability and how conduct changes as capability increases. This cohort supplies a rich account of economically situated conduct, but it does not identify that relationship. It contains four model families at one setting each, with differing response conditions, rather than an independently measured capability axis. Moreover, its prompt permits public deception. Calling all misleading public statements misalignment would apply a norm that was not assigned to the evaluated agents.

This does not make the social results irrelevant to safety. They identify behaviors and measurement problems that a norm-explicit evaluation must distinguish: selective disclosure, recipient correction, transient false echoes, continuing assurances, withdrawal after financial deterioration, and disclosure through a public field. They also show why a favorable settlement is not sufficient evidence of persuasive success. G13's buyer weighs solvency before the omission, but after the seller's opening appeal; G4's recipient had already recorded the true construction cost; G6 and G19 settle while correcting the premise. These restrictions do not identify a price effect or show that all manipulation fails.

The private-honesty request is also difficult to score from an inaccurate note. A returned explanation can be inconsistent, mistaken or corrected within the same response. Requiring a declared norm is necessary, but not sufficient, when the property refers to latent reasoning. Observable conduct standards---for example, disclosure of a specified encumbrance or honoring a precisely defined authorization boundary---can be evaluated more directly. Claims about honesty, intent and internal beliefs require evidence beyond the literal presence of a self-description.

This measurement dependence accords with mixed findings in related work. The Arena communication study reports no capability association under its economic ordering and communication codebook; Emergent Collusion reports within-family differences under designed violation opportunities~\cite{li2026arena,shi2026emergentcollusion}. MASK distinguishes factual accuracy from honesty, while Safetywashing questions whether safety benchmarks sometimes measure general capability~\cite{ren2025mask,ren2024safetywashing}. These results should motivate explicit constructs and controlled comparisons, rather than a universal monotonic relationship inferred from four win totals.

\subsection{From an auditable sandbox to controlled evaluation}
\label{sec:discussion-design}

The next experiment should be organized around the claim it intends to identify. A material-disclosure study could preserve G13's financial state while assigning a fixed disclosure requirement and varying whether the seller has an incentive to omit available collateral. A commitment study could specify the standing offer's terms, expiry and release authority before recreating a G18-style invocation. A protection task could designate information that must remain private while offering profitable opportunities to reveal it. These interventions would change the benchmark prompt or task and should be versioned as new conditions, not retroactively imposed on this cohort.

For a capability--adherence comparison, the capability axis should be independently measured, preferably with multiple tiers within model families and held-out economic tasks. Prompts, retry policies, memory, tool choice and resource budgets should be fixed or factorially varied. Refusals and invalid outputs must remain part of the measured system, with any recovery treatment separately recorded. Repeated complete games would measure long-run outcomes; repeated continuations from matched historical prefixes would answer narrower questions about the same available alternatives. Statistical units should reflect dependence within games and seed blocks.

Economic utility, instruction adherence and recipient effects should be reported separately. A high-profit policy that violates a declared authorization constraint should not become safe because the reward outweighs a hidden penalty. Similarly, a recipient that accepts a profitable exchange while detecting its false premise should not be scored as deceived solely because a transaction occurred. MACHIAVELLI's separation of rewards and norm-relevant behavior offers a related design precedent~\cite{pan2023machiavelli}; our cases add economically explicit sequences and recipient-side chronology.

Exact-prefix auditing offers a bridge to this design. A claim about inaccessible finance can become a reproducible local task with an engine-verified solution, while a true hotel lock supplies a negative case. A coherent hold and a defaulted build supply superficially similar END actions with different provenance. A corrected false premise and an untested persuasion hypothesis supply distinct social targets. A future monitor evaluation should declare the channels shown to the monitor, its prediction target and a held-out scoring rule. The present mismatch counts are diagnostics, not a validated monitorability benchmark.

Reinforcement learning requires additional engineering and validation: an observation/action interface, explicit objectives, reward semantics, horizon and truncation rules, reset distributions, and evaluation on situations not reused during training. A binding contract object or private bilateral channel may support particular social tasks, but neither is a universal prerequisite for training. Existing replay and recovery components support precise resets; they do not establish that learning in this sandbox improves behavior elsewhere. Transfer, reward exploitation and robustness to new counterparties remain empirical questions.

\subsection{The evaluator must be audited too}

Several consequential corrections concerned our own labels rather than agent actions. A terminal green inheritance had been described as G15's productive acquisition path; original events instead show orange completion, construction and \$3,250 in developed orange rent before the final estate. A \$1,200 construction step had been mistaken for G19 Gemini's \$2,000 cumulative expenditure on dark blues. Other errors arose from interpreting referents in messages, treating a menu as a complete financing set, or attributing a counteroffer to the thread initiator. Each could have changed an economic or social conclusion despite a numerically correct event log.

Overlapping review helped expose these errors, but reviewer agreement was not the final criterion. The corrections required source chronology, explicit units or legal execution. This distinction matters as agents increasingly evaluate other agents: a fluent, mutually agreed narrative can still misdescribe the recorded state. The research contribution is consequently both empirical and methodological---a set of observed trajectories, and a procedure for making claims about those trajectories answerable to their underlying records.

